\section{Related Work}

\textbf{Autonomic computing and self-healing.}
The self-healing concept originates in autonomic
computing~\cite{kephart2003vision}, which sought to embed
self-configuring, self-healing, self-optimising, and
self-protecting properties in distributed systems. Subsequent work
has applied the framework to security incidents at varying levels
of abstraction. Commercial implementations of self-healing
vulnerability remediation (Microsoft Defender automated
investigation; Tanium auto-remediation; IBM SOAR) are closed-source
and do not publish their decision rules or safety bounds.

\textbf{Pre-registered evaluation in security research.}
Pre-registration discipline is rare in security; the VulnPrio
sequence's earlier papers (Papers~1, 3--9) apply it consistently.
AutoHeal continues that discipline by locking the architecture's
parameters and hypotheses before evaluation.

\textbf{Integration with Papers 3--9.} AutoHeal synthesises:
HygieneBench (Paper~3)~\cite{paper3hygienebench} for the synthetic
telemetry substrate;
HygienePrio (Paper~4)~\cite{paper4hygieneprio} for scoring;
Paper~7's lag-1 calibration~\cite{paper7online} as the
calibration recipe at moderate capacity; Paper~9's
Theorem~1~\cite{paper9selftraj} (the closed-loop signal exhaustion
result) as the structural justification for the K=200 exclusion
from H1.

\textbf{Real CVE attribute distributions.} The frozen real
EPSS/KEV public corpus (\texttt{real\_data/processed/}) is the
same artifact released with the external-validity section of
Paper~4 (\S X)~\cite{paper4hygieneprio}.
AutoHeal inherits the corpus and its synthetic-vs-real distribution
gap as an inherited consequence: the absolute coverage numbers
under real distributions will differ from a purely-synthetic
evaluation by the same factor Paper~4 measured.

\textbf{Policy mandates.} CISA BOD~22-01~\cite{cisa2021bod2201}
and 23-01~\cite{cisa2022bod2301} require risk-based prioritization;
NIST SP~800-40 Rev.~4~\cite{nist2022sp80040r4} specifies the
documentation standard. AutoHeal's pre-registered triage rules
constitute the documentation form contemplated by these mandates,
extended with safety bounds that the mandates do not specify but
that production deployments require.

\textbf{What we are not the first to do.} Closed-loop patch
deployment is a long-standing operations practice. The novelty of
AutoHeal is not in the closed-loop concept but in the
\emph{pre-registered evaluation harness with locked safety bounds},
which is uncommon in both the academic and commercial literature.
